\documentclass[10pt,a4paper]{amsart}
\usepackage[margin=19mm]{geometry}
\usepackage{amsmath,amssymb,mathtools}
\usepackage[scr=rsfs,cal=euler]{mathalfa}
\usepackage{xfrac}
\usepackage[unicode,hidelinks,bookmarks=false]{hyperref}
\newcommand{\ie}{i.e.\ }
\newcommand{\cf}{cf.\ }
\newcommand{\resp}{respectively\ }
\newcommand{\Subsection}{\subsection}
\providecommand{\nobreakdash}{\nobreak\hbox{-}}
\setlength{\emergencystretch}{2em}
\multlinegap0pt
\usepackage{bm}
\newtheorem{proposition}{Proposition}
\usepackage{cleveref}
\crefname{proposition}{Proposition}{Propositions}
\newcommand{\Gmeas}{\bm{\Gamma}_{\bm R}}
\newcommand{\Gq}{\bm \Gamma_{\bm Q}}
\newcommand{\I}{{\mathbf{I}}}
\renewcommand{\L}{\mathbf{L}}
\renewcommand{\O}{{\mathbf{H}}}
\renewcommand{\S}{{\mathbf{S}}}
\newcommand{\Z}{{\mathbf{Z}}}
\DeclareMathOperator*{\argmax}{arg\,max}
\newcommand{\criteria}{\phi_\text{EIG}}
\newcommand{\logdet}{{\log\!\det}}
\renewcommand{\t}{{}^\top}
\title{Optimal sensor placement under model uncertainty in the weak-constraint 4D-Var framework}
\author{Alen Alexanderian, Hugo D\'iaz, Vishwas Rao, Arvind K. Saibaba}
\date{Source: arXiv:2502.00150v3 (CC BY 4.0)}
\begin{document}
\maketitle

\noindent\emph{Source and licence.} Optimal sensor placement under model uncertainty in the weak-constraint 4D-Var framework. Version arXiv:2502.00150v3 (CC BY 4.0), distributed by arXiv under the Creative Commons Attribution 4.0 International licence, http://creativecommons.org/licenses/by/4.0/, as recorded in the arXiv licence metadata for this exact version and shown on the abstract page https://arxiv.org/abs/2502.00150v3. Full licence notice: \texttt{licenses/arxiv-2502.00150-CC-BY-4.0.txt}.\par\smallskip

\noindent\textit{Editorial scope.} Counted unit: the article's proposition proposition:SCandWCoptCriteriaGap (source0.tex lines 921--939). The article's complete original proof of it is delivered with it (source0.tex lines 942--962, one \texttt{proof} environment of 21 lines). No mathematics is written, repaired or re-derived: the statement and the proof are the article's own lines, in the article's own notation, and no retained line is substituted or re-flowed.  Retained uncounted as the source-local context the proof needs: lines 874--888.  Nothing was omitted from any retained span, so the package carries no omission record.  No retained line contains a citation, so the wrapper carries no bibliography and no bibliography entry.  No retained line contains a figure environment, an image-inclusion command or drawing code, so no image is delivered and the package declares no asset dependency.  Editorial formatting only, in addition: an AMS \texttt{amsart} preamble that restates the article's own declarations and macros used by the retained lines, namely the environments \texttt{proposition} on separate counters, together with the standard packages those lines need.  The retained environments print in this document as Proposition 1 in order of appearance, whereas the article prints the same items with the numbering of its own full document.  The complete native source of the article as distributed in the arXiv e-print for this exact version is bundled unchanged as \texttt{sources/172550/source0.tex}.
\begin{proposition}\label{proposition:bound_SCandWCoptCriteria}
The following holds: 
\begin{align} \label{eq:gapEIG_SC_WC}
    0\leq \criteria  -   \criteria^\text{SC} 
      \leq \logdet \left( \I+ \bm Z
\begin{bmatrix}
 \bm 0 &    \\
   &   \Gq
\end{bmatrix}
\bm Z\t\right),
\end{align}
where $\Z := \Gmeas^{-\frac 12} \O \L^{-1}$.
Furthermore, as $\Gq \to \bm{0}$, the WC4D-Var EIG criterion converges to the SC4D-Var EIG criterion, i.e., $\criteria \to \criteria^\text{SC}$.

\end{proposition}

\begin{proposition}\label{proposition:SCandWCoptCriteriaGap}


Let \( \S^*_\text{WC} \) and \( \S^*_\text{SC} \) denote  optimal selection matrices for the WC4D-Var and SC4D-Var EIG criteria, namely:
\begin{align}
    \S^*_\text{WC} &\in \argmax_{\S} \criteria(\S), \quad
    \S^*_\text{SC} \in \argmax_{\S} \criteria^\text{SC}(\S).
\end{align}
Then, the gap between the criteria values for these optimal selections is bounded as:
\begin{align*}
    0\leq \criteria(\S^*_\text{WC}) - \criteria(\S^*_\text{SC}) \leq \logdet \left( \I + (\I \otimes \S^*_\text{WC})\bm{Z} 
    \begin{bmatrix}
    \bm{0} & \\
    & \Gq
    \end{bmatrix} 
    \bm{Z}^\top (\I \otimes \S^*_\text{WC})^\top \right).
\end{align*}

\end{proposition}

\begin{proof}
First, let \( \S \) be any selection matrix corresponding to the choice of \( k \) elements  out of \( n_s \). 
Then, the difference between the EIG criteria for the WC4D-Var and SC4D-Var frameworks satisfies the following bound:
\begin{align}\label{eq:gapEIG_SC_WC2}
    0 \leq \criteria(\S) - \criteria^\text{SC}(\S) \leq \logdet \left( \I + (\I \otimes \S)\bm{Z} 
    \begin{bmatrix}
    \bm{0} & \\
    & \Gq
    \end{bmatrix} 
    \bm{Z}^\top (\I \otimes \S)^\top \right),
\end{align}
where \( \bm{Z} := \Gmeas^{-\frac{1}{2}} \O \L^{-1} \).
The inequality \eqref{eq:gapEIG_SC_WC2} follows a similar approach to the proof of  \Cref{proposition:bound_SCandWCoptCriteria}. 
Now we turn to the actual result we wish to prove. The lower bound follows from the optimality of $\S^*_\text{WC}$.  For the upper bound, we decompose 
\[\begin{aligned}
    \criteria(\S^*_\text{WC}) - \criteria(\S^*_\text{SC}) = & \>  \underbrace{\criteria(\S^*_\text{WC}) - \criteria^\text{SC}(\S^*_\text{WC})}_{\equiv \alpha}
     + \underbrace{\criteria^\text{SC}(\S^*_\text{WC}) -  \criteria^\text{SC}(\S^*_\text{SC})}_{\equiv \beta}.
\end{aligned} \]
We have $\beta \le 0$ by  the optimality of $\S^*_\text{SC}$. The upper bound then follows from~\eqref{eq:gapEIG_SC_WC2} applied to $\S^*_\text{WC}$.
The uniform upper bound follows directly from Sylvester's law of inertia. 
\end{proof}

\end{document}
